% 西安交通大学本科毕业设计（论文）2026 Word 模板的 LaTeX 正文骨架。
% 编译：xelatex thesis.tex，至少两遍；真实文献数据库另按所用工具处理。
% 学校统一 A3 封面、任务书、考核评议书和答辩结果须使用正式文件另行并入。
% 此文件只提供 A4 论文内页；所有方括号占位文字须在提交前替换或删除。

\documentclass[UTF8,zihao=-4,twoside,openany,fontset=windows]{ctexbook}

\usepackage[
  a4paper,
  top=3.0cm,bottom=2.5cm,
  inner=2.6cm,outer=2.6cm,
  headheight=0.65cm,headsep=0.35cm,footskip=0.75cm
]{geometry}
\usepackage{fontspec}
\setmainfont{Times New Roman}
\setCJKmainfont{SimSun}
\usepackage{setspace}
\setstretch{1.2}
\setlength{\parindent}{2\ccwd}
\setlength{\parskip}{0pt}
\usepackage{fancyhdr}
\usepackage{tocloft}
\usepackage{graphicx}
\usepackage{caption}
\usepackage{booktabs,tabularx}
\usepackage{amsmath,amssymb}
\usepackage[numbers,super,sort&compress]{natbib}

% 章、节、小节分别为三号、小三号、四号；章节编号为 1 / 1.1 / 1.1.1。
\ctexset{
  contentsname={目\quad 录},
  chapter={
    name={,},number=\arabic{chapter},aftername=\quad,
    format=\normalfont\songti\zihao{3}\centering,
    beforeskip=2\baselineskip,afterskip=1\baselineskip
  },
  section={
    format=\normalfont\songti\zihao{-3}\raggedright,
    aftername=\quad,indent=0pt,
    beforeskip=1\baselineskip,afterskip=.5\baselineskip
  },
  subsection={
    format=\normalfont\songti\zihao{4}\raggedright,
    aftername=\quad,indent=2\ccwd,
    beforeskip=.5\baselineskip,afterskip=0pt
  },
  appendix/name={附录,},appendix/number=\Alph{chapter}
}
\setcounter{secnumdepth}{2}
\setcounter{tocdepth}{2}

% 目录：三号宋体标题；小四号宋体条目；逐级缩进一字符，右侧页码用点线连接。
\renewcommand{\cfttoctitlefont}{\hfill\normalfont\songti\zihao{3}}
\renewcommand{\cftaftertoctitle}{\hfill}
\setlength{\cftbeforetoctitleskip}{2\baselineskip}
\setlength{\cftaftertoctitleskip}{1\baselineskip}
\setlength{\cftchapindent}{0pt}
\setlength{\cftsecindent}{1\ccwd}
\setlength{\cftsubsecindent}{2\ccwd}
\setlength{\cftchapnumwidth}{2.5em}
\setlength{\cftsecnumwidth}{3.5em}
\setlength{\cftsubsecnumwidth}{4.5em}
\renewcommand{\cftchapfont}{\normalfont\songti\zihao{-4}}
\renewcommand{\cftsecfont}{\normalfont\songti\zihao{-4}}
\renewcommand{\cftsubsecfont}{\normalfont\songti\zihao{-4}}
\renewcommand{\cftchappagefont}{\normalfont\zihao{-4}}
\renewcommand{\cftsecpagefont}{\normalfont\zihao{-4}}
\renewcommand{\cftsubsecpagefont}{\normalfont\zihao{-4}}
\renewcommand{\cftchapleader}{\cftdotfill{\cftdotsep}}

% 奇数页显示对应章名，偶数页显示校名；页码在页脚外侧。
\renewcommand{\chaptermark}[1]{\markboth{\thechapter\quad#1}{}}
\newcommand{\xjtuheadfoot}{%
  \fancyhf{}%
  \fancyhead[CO]{\zihao{5}\nouppercase{\leftmark}}%
  \fancyhead[CE]{\zihao{5}西安交通大学本科毕业设计（论文）}%
  \fancyfoot[LE,RO]{\zihao{5}\thepage}%
  \renewcommand{\headrulewidth}{0.5pt}%
  \renewcommand{\headrule}{%
    \vbox{\hrule width\headwidth height0.5pt\vskip1.5pt
          \hrule width\headwidth height0.5pt}}%
}
\pagestyle{fancy}
\xjtuheadfoot
\fancypagestyle{plain}{\xjtuheadfoot}

% 图表题注五号居中；图题在下，表题在上。
\DeclareCaptionFont{xjtuwu}{\zihao{5}}
\DeclareCaptionLabelFormat{xjtu}{#1#2}
\DeclareCaptionLabelFormat{xjtuappendix}{#2}
\captionsetup{font=xjtuwu,justification=centering,labelformat=xjtu,labelsep=space}
\renewcommand{\thefigure}{\thechapter-\arabic{figure}}
\renewcommand{\thetable}{\thechapter-\arabic{table}}
\renewcommand{\theequation}{\thechapter-\arabic{equation}}
\setlength{\heavyrulewidth}{1.5pt}
\setlength{\lightrulewidth}{1pt}
\makeatletter
\renewcommand{\tagform@}[1]{\maketag@@@{（#1）}}
\makeatother

\begin{document}

% 正文前页码为大写罗马数字。中文、英文摘要各占新页。
\pagenumbering{Roman}
\chapter*{摘\quad 要}
\markboth{摘\quad 要}{}
\addcontentsline{toc}{chapter}{摘\quad 要}
[填写中文摘要：目的、方法、本人工作、主要结果和结论。]

\par\noindent\textbf{关\quad 键\quad 词：}[词一]；[词二]；[词三]

\clearpage
\chapter*{ABSTRACT}
\markboth{ABSTRACT}{}
\addcontentsline{toc}{chapter}{ABSTRACT}
[Write the English abstract corresponding to the Chinese abstract.]

\par\noindent\textbf{KEY WORDS:} [term one]; [term two]; [term three]

% 符号较多时，在目录前启用以下部分；符号少时删除此注释块。
% \clearpage
% \chapter*{主要符号表}
% \markboth{主要符号表}{}
% \addcontentsline{toc}{chapter}{主要符号表}
% \noindent [按拉丁大写、小写、希腊大写、小写排列，并注明单位及缩略词全称。]

% 中文目录须从奇数页起排。原 Word 文件的页序与装订条款不同；此处按装订条款。
\cleardoublepage
\tableofcontents

% 正文第一章从奇数页起排，从阿拉伯数字 1 开始。
\cleardoublepage
\pagenumbering{arabic}
\chapter{绪论}
\section{研究背景}
[填写真实研究背景。]
\subsection{研究问题}
[填写研究问题、已有工作及本文任务。]

% 根据真实论文需要添加章节；不要保留模板示例章节数。
% 图表示例：先在正文中引用“图\ref{fig:example}”，再插入图。
% \begin{figure}[htbp]
%   \centering
%   \includegraphics[width=6.67cm]{figures/example.pdf}
%   \caption{真实图题}\label{fig:example}
% \end{figure}
%
% 三线表示例：表题必须在表格上方，表宽与正文版心一致。
% \begin{table}[htbp]
%   \centering
%   \caption{真实表题}\label{tab:example}
%   \zihao{5}
%   \begin{tabularx}{\textwidth}{lXX}
%     \toprule
%     项目 & 数据一 & 数据二 \\
%     \midrule
%     [实际项目] & [实际值] & [实际值] \\
%     \bottomrule
%   \end{tabularx}
% \end{table}

\chapter{结论与展望}
[归纳本人研究结果、适用范围、不足与后续建议。]

\chapter*{致\quad 谢}
\markboth{致\quad 谢}{}
\addcontentsline{toc}{chapter}{致\quad 谢}
[填写真实致谢；不超过 1000 汉字。]

\chapter*{参考文献}
\markboth{参考文献}{}
\addcontentsline{toc}{chapter}{参考文献}
% 填入真实文献。使用 BibTeX/Biber 时，替换此标题及以下区域，
% 确保自动生成的参考文献标题不会与本标题重复，并加入目录。

% 如需附录，取消以下示例的注释并填入真实内容；页码继续正文编号。
% \appendix
% \renewcommand{\thefigure}{附录\Alph{chapter}-图\arabic{figure}}
% \renewcommand{\thetable}{附录\Alph{chapter}-表\arabic{table}}
% \renewcommand{\theequation}{附录\Alph{chapter}-\arabic{equation}}
% \captionsetup[figure]{labelformat=xjtuappendix}
% \captionsetup[table]{labelformat=xjtuappendix}
% \chapter{外文原文}
% [填写真实附录内容。]

\end{document}
